\paragraph{Question and expectation.} Can the surrogate find better designs by itself? From the six Gaussian processes
the smallest reserve was predicted for every conceivable design, and one standard deviation was subtracted (cautious
estimate). The four designs with the highest value were fully computed. All four have 24 coils, a coil--plasma
distance of \SI{0.21}{m}, short coils (\SIrange{3.9}{4.1}{m} per coil) and the thickest outer wall; they differ
mainly in number and stiffness of the connectors. The cautious prediction was 0.04--0.07. Expected (event log
F-142): 1--3 of the 4 admissible, stress $\pm30$~\% off the prediction, none clearly more robust than 0169.

\begin{table}[H]\centering\scriptsize\begin{tabular}{@{}lrrrrrrrrl@{}}\toprule
Proposal & $E_\text{c}$ & c./pair & mass & $\sigma_{99}$ & $\sigma_{99}$ fine & strain & $B\cdot n$ & \multicolumn{2}{c}{robustness} \\
 & GPa & & kt & MPa & MPa & \% & ‰ & predicted & computed (bottleneck) \\\midrule
V0 & \num{55} & 13 & \num{37.8} & \num{306} & \num{385} & \num{0.16} & \num{1.52} & 0.07 & 0.23 (package space) \\
V1 & \num{77} & 13 & \num{33.1} & \num{220} & \num{366} & \num{0.12} & \num{1.87} & 0.06 & 0.31 (package space) \\
V2 & \num{81} & 9 & \num{26.4} & \num{851} & -- & \num{0.54} & \num{2.19} & 0.06 & -0.79 (winding-pack strain), inadmissible \\
V3 & \num{176} & 11 & \num{29.1} & \num{272} & \num{394} & \num{0.13} & \num{1.96} & 0.04 & 0.35 (field error B·n) \\
\bottomrule\end{tabular}
\caption{Surrogate proposals. $E_\text{c}$ connector
stiffness, c./pair connectors per coil pair, ``fine'' = 1.5 times as many stations. Robustness = smallest reserve over
all eight limits; ``predicted'' = prediction minus one standard deviation.}
\end{table}

\paragraph{Result.} Three of the four proposals are admissible, and the computed robustness (0.23--0.35) is far above
the cautious prediction. V3 (11 stiff connectors per pair, \SI{176}{GPa}) reaches 0.35 and is, on the standard mesh,
the most robust of all designs computed so far; its bottleneck is the field error, not stress. V2, with only 9
connectors per pair, fails clearly: $\sigma_{99}$ \SI{851}{MPa}, winding-pack strain 0.54~\%, deformation
\SI{177}{mm}. The surrogate did not see this outlier coming.

\paragraph{Mesh check.} As for the doe1 candidates, $\sigma_{99}$ rises markedly with 1.5 times as many stations: V3
\SI{272}{} $\to$ \SI{394}{MPa} ($+45$~\%), V1 \SI{220}{} $\to$ \SI{366}{MPa} ($+66$~\%), V0 \SI{306}{} $\to$
\SI{385}{MPa} ($+26$~\%). All remain admissible, but V3 drops to a robustness of 0.21 because steel now limits. 0169
(\SI{319}{MPa} on the fine mesh, robustness 0.29) remains the most robust design.

\paragraph{Interpretation.} The surrogate is useful for pre-selection: it reliably finds admissible regions (it
predicts package, nuclear heating and field error well). For stress it is not enough; a single proposal can be far
off, and it does not see the mesh dependence of the peaks at all. Notably, 0169 remains the only design whose stress
hardly rises on the finer mesh. In 0169 the connectors merge into almost continuous plates (section~\ref{sec:3d}). A
conjecture, not yet tested: fewer free connector ends mean fewer stiffness jumps and hence peaks that do not grow
with the mesh.
